\documentclass[10pt,a4paper]{amsart}
\usepackage[margin=19mm]{geometry}
\usepackage{amsmath,amssymb,mathtools}
\usepackage[scr=rsfs,cal=euler]{mathalfa}
\usepackage{xfrac}
\usepackage[unicode,hidelinks,bookmarks=false]{hyperref}
\newcommand{\ie}{i.e.\ }
\newcommand{\cf}{cf.\ }
\newcommand{\resp}{respectively\ }
\newcommand{\Subsection}{\subsection}
\providecommand{\nobreakdash}{\nobreak\hbox{-}}
\setlength{\emergencystretch}{2em}
\multlinegap0pt
\usepackage{bm}
\usepackage{bbm}
\usepackage{dsfont}
\usepackage{xspace}
\usepackage{cancel}
\usepackage{mathtools}
\usepackage{amsthm}
\makeatletter
\@ifundefined{theorem}{\newtheorem{theorem}{Theorem}}{}
\makeatother
\makeatletter
\expandafter\ifx\csname proof\endcsname\relax
\renewenvironment{proof}{{\noindent\bfseries Proof.}}{\qed}
\fi
\makeatother
\title[Combinatorial and analytic aspects of independence polynomia]{Combinatorial and analytic aspects of independence polynomials of zero divisor graphs}
\author{Bilal Ahmad Rather}
\date{Source: arXiv:2606.04789v1 (CC BY 4.0)}
\begin{document}
\maketitle

\noindent\emph{Source and licence.} Combinatorial and analytic aspects of independence polynomials of zero divisor graphs. Version arXiv:2606.04789v1 (CC BY 4.0), distributed under the Creative Commons Attribution 4.0 International licence, as recorded in the arXiv licence metadata for this exact version and shown on the abstract page https://arxiv.org/abs/2606.04789v1. Full rights notice: licenses/arxiv-2606.04789-CC-BY-4.0.txt.\par\smallskip

\noindent\textit{Editorial scope.} Counted unit: the Theorem whose source label is \texttt{None} in the article (source0.tex lines 722--727, source label \texttt{None}), delivered together with the article's own complete original proof of it (source0.tex lines 729--995, one \texttt{proof} environment of 267 lines). No mathematics is written, repaired or re-derived: statement and proof are the article's own text line for line. Required context retained uncounted: none. Every definition, display and label of those lines is retained unchanged. Omitted from this package and not redistributed: 1--721, 728--728, 996--1375. Nothing inside a retained excerpt is omitted or altered: every retained body is delivered exactly as the article's lines are written, so no omission receipt of any kind accompanies this package. Citations: the retained excerpts carry no citation command at all, so no bibliography entry of the article is transcribed and no reference is redistributed. Reference adaptation: none was needed and none was made; every reference inside the delivered unit resolves inside it, so no unresolved reference or citation is produced. Printed slips kept literal: no printed word, symbol, hypothesis or number of the counted statement or of its proof was altered. Layout and class: the article's own class is \texttt{article} with packages including graphicx, graphics, amsthm, amsfonts, amsbsy, amssymb, amsmath, cite, array, arydshln, hyperref, float, tikz, standalone; the wrapper is the lane's pinned \texttt{amsart} editorial wrapper at 10pt on A4, which restates the theorem-like environments the retained text needs and adds the article's own macro definitions that the retained text needs (none), with extra wrapper packages (bm, bbm, dsfont, xspace, cancel, mathtools). The delivered numbering is the unit's own. Assets: no retained span contains a figure environment, a graphics-inclusion command, a PDF-page inclusion, a table or a TikZ picture; no image file is copied into this package, the package declares no asset dependency and no image file is redistributed with it. Licence: version arXiv:2606.04789v1 of this article is distributed under the Creative Commons Attribution 4.0 International licence, as recorded in the arXiv licence metadata for this exact version and shown on the abstract page https://arxiv.org/abs/2606.04789v1. A full rights notice is delivered at licenses/arxiv-2606.04789-CC-BY-4.0.txt. Characters: every retained excerpt is pure ASCII, so the delivered engine, XeLaTeX with its default Latin Modern text font, reports no missing character.
\begin{theorem}
	Let $I(x)$ be the independence polynomial of $\varGamma(\mathbb{Z}_{p^{2}q})$, where $p<q$ are primes. Then  all zeros of $I(x)$ different from $x=-1$ lie in the annular region
	$$
	\left\{x\in\mathbb{C}:\frac{1}{2}<|1+x|<2\right\}.
	$$
\end{theorem}

\begin{proof}
	Since the independence polynomial of $\varGamma(\mathbb{Z}_{p^{2}q})$ is
	$$
	I(x)
	=
	(1+x)^{(p-1)(p+q-1)}
	+
	(1+x)^{p(q-1)}
	-
	(1+x)^{(p-1)(q-1)}
	+
	(p-1)x(1+x)^{p(p-1)}.
	$$
	With 
	$
	y=1+x,
	$ the above polynomial becomes
	$$
	I(y-1)
	=
	y^{(p-1)(p+q-1)}
	+
	y^{p(q-1)}
	-
	y^{(p-1)(q-1)}
	+
	(p-1)(y-1)y^{p(p-1)}.
	$$
	Since $p(p-1)>0$, it follows immediately that $y=0$, that is $x=-1$, is a zero of $I(x)$. We are interested only in the other zeros, so we assume $y\neq 0$.  The polynomial can be written as $
	I(y-1)
	=
	y^{p(p-1)}Q(y),
	$
	where
	$$
	Q(y)
	=
	y^{A}+y^{B}-y^{C}+(p-1)y-(p-1),
	$$
	with
	$
	A=(p-1)(q-1),
	B=p(q-p),$ and $
	C=(p-1)(q-p-1).
	$
	Thus the zeros of $I(x)$ different from $x=-1$ correspond exactly to the non-zero zeros of $Q(y)$. It is therefore enough to prove that every zero $y$ of $Q(y)$ satisfies
	$
	\tfrac{1}{2}<|y|<2.
	$
	We shall prove separately that $Q(y)\neq 0$ for $|y|\leq \tfrac12$ and for $|y|\geq 2$.
	 First suppose that $|y|\leq \tfrac12$, and we will show that $Q(y)\neq 0$. With $r=|y|$, if $Q(y)=0$, then
	$$
	(p-1)
	=
	y^{A}+y^{B}-y^{C}+(p-1)y.
	$$
	Now, with absolute values, we have
	$$
	p-1
	\leq
	r^{A}+r^{B}+r^{C}+(p-1)r.
	$$
	Next, we check that the right-hand side of above inequality is always strictly smaller than $p-1$.
	 If $(p,q)=(2,3)$, then
	$
	Q(y)=2y^{2}+y-2.
	$
	If $|y|\leq \tfrac12$ and $Q(y)=0$, then
	$$
	2=|2y^{2}+y|\leq 2|y|^{2}+|y|\leq 2\left(\frac12\right)^2+\frac12=1,
	$$
	which is impossible. Next, let $p=2$ and $q\geq 5$. Then
	$
	Q(y)=y^{q-1}+y^{2q-4}-y^{q-3}+y-1.
	$
	If $Q(y)=0$, then
	$
	1
	\leq
	r^{q-1}+r^{2q-4}+r^{q-3}+r.
	$
	Since $q\geq 5$ and $r\leq \frac12$, we obtain
	$$
	1
	\leq
	\left(\frac12\right)^4
	+
	\left(\frac12\right)^6
	+
	\left(\frac12\right)^2
	+
	\frac12
	=
	\frac{53}{64}<1,
	$$
	a contradiction. For $p=3$, since $p<q$ and both $3$ and $q$ are primes, so we have $q\geq 5$. Hence,
	$$
	A=2(q-1)\geq 8,\qquad
	B=3(q-3)\geq 6,\qquad
	C=2(q-4)\geq 2.
	$$
	Thus, if $Q(y)=0$ and $|y|\leq \frac12$, then
	$$
	2
	\leq
	\left(\frac12\right)^8
	+
	\left(\frac12\right)^6
	+
	\left(\frac12\right)^2
	+
	2\cdot\frac12
	<2,
	$$
	which is again impossible. For $p\geq 5$, since $p$ and $q$ are odd primes with $p<q$, so we have $q\geq p+2$. Hence, 
	$
	C=(p-1)(q-p-1)\geq p-1\geq 4.
	$
	Also $A\geq 4$ and $B\geq 4$. Therefore, we derive
	$$
	r^{A}+r^{B}+r^{C}+(p-1)r
	\leq
	3\left(\frac12\right)^4+\frac{p-1}{2}
	=
	\frac{3}{16}+\frac{p-1}{2}.
	$$
	Since $p\geq 5$, we have $p-1\geq 4$, and so
	$
	\frac{3}{16}+\frac{p-1}{2}<p-1.
	$
	This contradicts the inequality obtained from $Q(y)=0$. Hence, $Q(y)$ has no zeros in $|y|\leq \frac12$.
	
	It remains to show that $Q(y)$ has no zeros in $|y|\geq 2$. Let $r=|y|\geq 2$. We use the relations
	$
	A-C=p(p-1), 
	B-C=q-1,
	$
	and
	$
	A-B=p(p-1)-(q-1).
	$
	We have to consider three cases. If $A>B$. Then $A-B\geq 1$. If $Q(y)=0$, then
	$
	y^A=-y^B+y^C-(p-1)y+(p-1).
	$
	With absolute values, we get
	$$
	1
	\leq
	r^{B-A}+r^{C-A}+(p-1)r^{1-A}+(p-1)r^{-A}.
	$$
	Since $A-C=p(p-1)$, and $r\geq 2$, this implies that
	$$
	1
	\leq
	2^{-(A-B)}
	+
	2^{-p(p-1)}
	+
	(p-1)2^{-(A-1)}
	+
	(p-1)2^{-A}.
	$$
	For case $A>B$, we necessarily have $p\geq 3$. Since $q\geq p+2$, we obtain
	$
	A=(p-1)(q-1)\geq (p-1)(p+1),
	$
	and hence $A-1\geq p(p-1)$. Thus, we obtain
	$$
	1
	\leq
	\frac12
	+
	2^{-p(p-1)}
	+
	(p-1)2^{-p(p-1)}
	+
	(p-1)2^{-p(p-1)-1},
	$$
	or equivalently,
	$$
	1
	\leq
	\frac12+
	\frac{1+\frac32(p-1)}{2^{p(p-1)}}.
	$$
	But $p\geq 3$, so $p(p-1)\geq 6$, and hence
	$$
	\frac{1+\frac32(p-1)}{2^{p(p-1)}}<\frac12.
	$$
	Therefore, the right-hand side is strictly less than $1$, a contradiction. Hence $Q(y)\neq 0$ for $|y|\geq 2$ in the case $A>B$.
	 Now, for $A=B$, we get 
	$
	Q(y)=2y^A-y^C+(p-1)y-(p-1).
	$
	If $Q(y)=0$, then
	$
	2r^A\leq r^C+(p-1)r+(p-1), 
	$ and thereby we have
	$$
	2\leq r^{C-A}+(p-1)r^{1-A}+(p-1)r^{-A}.
	$$
	Since $A-C=p(p-1)$ and $r\geq 2$, the right-hand side of above inequality  is at most
	$$
	2^{-p(p-1)}+(p-1)2^{1-A}+(p-1)2^{-A},
	$$
	which is strictly less than $2$. Clearly, for $(p,q)=(2,3)$, it is at most
	$
	\frac14+\frac12+\frac14=1,
	$
	and for $p\geq 3$ it is even smaller, since $A$ and $p(p-1)$ are larger. Thus, we get a contradiction. Therefore $Q(y)\neq 0$ for $|y|\geq 2$ in the case $A=B$. Finally, if $B>A$, then
	$
	B-A=(q-1)-p(p-1)>0.
	$
	Since $q$ is odd and $p(p-1)$ is even, the positive integer $B-A$ is in fact at least $2$. If $Q(y)=0$, then
	$$
	y^B=-y^A+y^C-(p-1)y+(p-1),
	$$
	and hence
	$$
	1
	\leq
	r^{A-B}+r^{C-B}+(p-1)r^{1-B}+(p-1)r^{-B}.
	$$
	With $B-C=q-1$, we obtain
	$$
	1
	\leq
	2^{-(B-A)}
	+
	2^{-(q-1)}
	+
	(p-1)2^{1-B}
	+
	(p-1)2^{-B}.
	$$
	Since $B-A\geq 2$ and $q-1\geq 4$, the first two terms are bounded by $\tfrac14$ and $\tfrac1{16}$, respectively. Also, we note that 
	$
	(p-1)2^{1-B}+(p-1)2^{-B}
	=
	\tfrac{3(p-1)}{2^B}.
	$
	If $p=2$, then $q\geq 5$, so $B=2(q-2)\geq 6$, and therefore
	$
	\tfrac{3(p-1)}{2^B}\leq \tfrac{3}{64}.
	$
	If $p\geq 3$, then $B>A$ implies that $q-1>p(p-1)$, and so $B=p(q-p)$ is sufficiently large. In particular,
	$
	\tfrac{3(p-1)}{2^B}<\tfrac14.
	$
	Thus,  in all cases we get
	$
	1
	<
	\tfrac14+\tfrac1{16}+\tfrac14
	<1,
	$
	which is absurd. Therefore, $Q(y)$ has no zeros with $|y|\geq 2$ in the case $B>A$.
	Thus, by combining the two parts, every zero $y$ of $Q(y)$ satisfies
	$
	\tfrac12<|y|<2.
	$
	Since the zeros of $I(x)$ different from $x=-1$ correspond to the zeros of $Q(y)$ under the change of variable $y=1+x$, we conclude that every zero $x\neq -1$ of $I(x)$ satisfies
	$
	\tfrac12<|1+x|<2.
	$
\end{proof} 

\end{document}
